% Part 3 of Day 3: Features for time series.
%
% Topics of the syllabus: statistical features and spectral features; why
% signal processing before the network reduces the model size. Exercise:
% compare the input size of raw windows and of spectral features.
%
% The window of the motion lab has 2 s at 50 Hz with 3 axes: 100 samples for
% each axis. The feature block of the kit lab gives 21 features for each
% axis with an FFT length of 32: 63 features in total.

\theorypart{Features for time series}%
  {compare the input size of raw windows and of spectral features}

% --------------------------------------------------------------- two inputs
\begin{frame}{Two ways to give a window to a model}
  \begingroup\centering
  \begin{tikzpicture}[font=\footnotesize, >={Stealth[length=5pt]},
      b/.style={draw=coolgray, rounded corners=3pt, fill=boxgray,
                minimum width=2.1cm, minimum height=0.9cm, align=center}]
    \node[anchor=west, text=coolgray] at (-1.05,0.9) {\textbf{Raw window}};
    \node[b] (a1) at (0,0)   {Window\\300 values};
    \node[b, minimum width=3.4cm] (a3) at (5.9,0) {Network with\\300 inputs};
    \node[b] (a4) at (8.9,0) [minimum width=1.4cm] {Class};
    \draw[->] (a1) -- (a3); \draw[->] (a3) -- (a4);
    \node[anchor=west, text=coolgray] at (-1.05,-1.2) {\textbf{Features}};
    \node[b] (b1) at (0,-2.1)   {Window\\300 values};
    \node[b, draw=cardinalred] (b2) at (2.9,-2.1) {Signal\\processing};
    \node[b] (b3) at (5.9,-2.1) {Network with\\63 inputs};
    \node[b] (b4) at (8.9,-2.1) [minimum width=1.4cm] {Class};
    \draw[->] (b1) -- (b2);
    \draw[->] (b2) -- (b3);
    \draw[->] (b3) -- (b4);
  \end{tikzpicture}\par\endgroup
  \vskip 0.5em
  \begin{itemize}
    \item A feature is one number that describes a property of the window:
          how strong the motion is, how fast it is.
    \item Signal processing calculates the features with fixed formulas. It
          has no trained weights.
    \item The network then learns only the step from the features to the
          class.
  \end{itemize}
  \keypoint{The question of this part: which input gives the smaller model
    and the smaller arena for the same task?}
\end{frame}

% ------------------------------------------------------- statistical features
\begin{frame}{Statistical features in the time domain}
  \small
  \renewcommand{\arraystretch}{1.3}
  \begin{tabular}{@{}L{0.20\textwidth}L{0.34\textwidth}L{0.36\textwidth}@{}}
    \toprule
    \textbf{Feature} & \textbf{What it measures}
    & \textbf{In the motion classes} \\
    \midrule
    Mean      & The constant part of the signal
              & Gravity: the angle of the kit \\
    RMS       & The energy of the signal
              & Near 0 for ``idle'', large for a strong motion \\
    Minimum, maximum & The range of the values & The size of the motion \\
    Skewness  & If the values are symmetric around the mean
              & A motion with a fast and a slow direction \\
    Kurtosis  & If the signal has short, large peaks & A shock \\
    \bottomrule
  \end{tabular}
  \vskip 0.7em
  \begin{itemize}
    \item Each feature is one number for each axis. The order of the
          samples in the window has no effect on these five features.
    \item You calculate them with one loop over the samples.
  \end{itemize}
  \source{Feature list: \emph{Machine Learning Systems}, hardware kits,
    ``DSP Spectral Features''.}
\end{frame}

\begin{frame}{RMS: the energy of the motion}
  \begin{columns}[c]
    \begin{column}{0.50\textwidth}
      \[
        x_{\text{RMS}} = \sqrt{\frac{1}{n}\left(x_1^2 + x_2^2 + \dots +
          x_n^2\right)}
      \]
      \begin{itemize}
        \item RMS: root mean square.
        \item Subtract the mean of the window first. If not, gravity
              decides the value of the z axis.
        \item After this step, the RMS is equal to the standard
              deviation.
      \end{itemize}
    \end{column}
    \begin{column}{0.46\textwidth}
      \centering
      \begin{tikzpicture}[font=\scriptsize, x=1.85cm, y=0.9cm]
        \draw[->, coolgray] (0,0) -- (2.1,0) node[right] {s};
        \draw[->, coolgray] (0,-1.3) -- (0,1.4);
        \draw[darkcyan, line width=0.9pt]
          plot[domain=0:2, samples=100, smooth] (\x,{sin(360*1.5*\x)});
        \draw[cardinalred, dashed] (0,0.707) -- (2,0.707);
        \draw[cardinalred, dashed] (0,-0.707) -- (2,-0.707);
        \node[cardinalred, anchor=south east] at (2,0.707) {RMS};
        \node[anchor=east] at (0,1) {1};
        \node[anchor=east] at (0,-1) {$-1$};
      \end{tikzpicture}
      \vskip 0.2em
      {\footnotesize A sine wave with the amplitude 1 has an RMS of
       $1 / \sqrt{2} = 0.707$.\par}
    \end{column}
  \end{columns}
  \vskip 0.5em
  \keypoint{The RMS of three axes gives three numbers. They separate
    ``idle'' from motion, and a horizontal motion from a vertical motion.}
\end{frame}

\begin{frame}{Skewness and kurtosis: the form of the distribution}
  \begingroup\centering
  \begin{tikzpicture}[font=\scriptsize, x=0.55cm, y=1.6cm]
    % three distributions of the sample values of a window
    \foreach \o/\t in {0/symmetric, 6.5/skewed, 13/peaks} {
      \draw[->, coolgray] (\o-2.6,0) -- (\o+2.8,0);
      \node[anchor=north] at (\o,-0.08) {\textbf{\t}};
    }
    \draw[darkcyan, line width=0.9pt]
      plot[domain=-2.5:2.5, samples=60, smooth] (\x,{exp(-\x*\x)});
    \draw[cardinalred, line width=0.9pt]
      plot[domain=0:5, samples=60, smooth]
      (\x+4.0,{2.72*0.9*\x*exp(-1.8*\x) + 0.0});
    \draw[treegreen, line width=0.9pt]
      plot[domain=-2.5:2.5, samples=80, smooth]
      (\x+13,{0.8*exp(-12*\x*\x) + 0.2*exp(-0.5*\x*\x)});
    \node at (0,-0.45) {skewness 0};
    \node at (6.5,-0.45) {skewness above 0};
    \node at (13,-0.45) {high kurtosis};
  \end{tikzpicture}\par\endgroup
  \vskip 0.4em
  \begin{itemize}
    \item Each curve shows how often each value occurs in a window.
    \item \textbf{Skewness} measures the symmetry. A motion that goes fast
          in one direction and slowly back gives a skewed distribution.
    \item \textbf{Kurtosis} measures the weight of the rare, large values. A
          signal that is quiet with some shocks has a high kurtosis.
    \item The RMS says how large the signal is. These two features say
          which form it has.
  \end{itemize}
  \source{The curves are drawings. Definitions: \emph{Machine Learning
    Systems}, hardware kits, ``DSP Spectral Features''.}
\end{frame}

\begin{frame}{The statistics do not see the speed}
  \begingroup\centering
  \begin{tikzpicture}[font=\scriptsize, x=4.4cm, y=0.8cm]
    \foreach \o/\f/\t in {0/1/slow motion: 1 Hz, 1.2/4/fast motion: 4 Hz} {
      \draw[->, coolgray] (\o,0) -- (\o+1.08,0);
      \draw[->, coolgray] (\o,-1.2) -- (\o,1.3);
      \draw[darkcyan, line width=0.9pt]
        plot[domain=0:1, samples=120, smooth] (\x+\o,{sin(360*\f*\x)});
      \node[anchor=south] at (\o+0.5,1.25) {\textbf{\t}};
      \node at (\o+0.5,-1.6) {mean 0, RMS 0.707};
    }
  \end{tikzpicture}\par\endgroup
  \vskip 0.4em
  \begin{itemize}
    \item The two signals have the same mean, the same RMS, the same
          minimum, and the same maximum.
    \item The statistical features are equal. A model with these features
          cannot separate the two motions.
    \item The difference is the frequency. You need a feature that
          measures it.
  \end{itemize}
  \vskip 0.2em
  \keypoint{The time domain says how large the signal is. The frequency
    domain says how fast it changes.}
\end{frame}

% ---------------------------------------------------------- frequency domain
\begin{frame}{The frequency domain}
  \begin{columns}[c]
    \begin{column}{0.50\textwidth}
      \centering
      \begin{tikzpicture}[font=\scriptsize, x=4.6cm, y=0.75cm]
        \draw[->, coolgray] (0,0) -- (1.08,0) node[right] {s};
        \draw[->, coolgray] (0,-1.7) -- (0,1.8);
        \draw[darkcyan, line width=0.9pt]
          plot[domain=0:1, samples=160, smooth]
          (\x,{sin(360*2*\x) + 0.5*sin(360*6*\x)});
        \node at (0.5,-2.1) {time domain};
      \end{tikzpicture}
    \end{column}
    \begin{column}{0.46\textwidth}
      \centering
      \begin{tikzpicture}[font=\scriptsize, x=0.42cm, y=1.5cm]
        \draw[->, coolgray] (0,0) -- (9.4,0) node[right] {Hz};
        \draw[->, coolgray] (0,0) -- (0,1.3);
        \foreach \f in {0,2,4,6,8} { \node[below, coolgray] at (\f,0) {\f}; }
        \draw[cardinalred, line width=3pt] (2,0) -- (2,1);
        \draw[cardinalred, line width=3pt] (6,0) -- (6,0.5);
        \node[anchor=west] at (2.3,1.0) {amplitude 1};
        \node[anchor=west] at (6.3,0.5) {0.5};
        \node at (4.5,-0.52) {frequency domain};
      \end{tikzpicture}
    \end{column}
  \end{columns}
  \vskip 0.3em
  \begin{itemize}
    \item Each signal is a sum of sine waves with different frequencies.
    \item The signal on the left is the sum of two sine waves: 2 Hz with
          the amplitude 1, and 6 Hz with the amplitude 0.5.
    \item The spectrum on the right shows the same signal: one bar for each
          frequency.
    \item The Fourier transform calculates the spectrum from the samples.
          The FFT is a fast method for it.
  \end{itemize}
\end{frame}

\begin{frame}{The FFT of a window}
  \begin{columns}[T]
    \begin{column}{0.50\textwidth}
      \begin{itemize}
        \item The FFT takes $N$ samples. $N$ is the FFT length, a power
              of two.
        \item It gives $N / 2 + 1$ values. Each value is one frequency bin.
        \item The width of one bin is $f_s / N$.
        \item The last bin is at $f_s / 2$, the Nyquist frequency of Day 2.
      \end{itemize}
    \end{column}
    \begin{column}{0.46\textwidth}
      \begin{block}{The motion lab}
        $f_s = 50$ Hz and $N = 32$:
        \begin{itemize}
          \item 17 bins, from 0 to 25 Hz
          \item bin width $50 / 32 = 1.56$ Hz
        \end{itemize}
      \end{block}
    \end{column}
  \end{columns}
  \vskip 0.5em
  \small
  \renewcommand{\arraystretch}{1.25}
  \begin{tabular}{@{}lrrl@{}}
    \toprule
    \textbf{FFT length} & \textbf{Bins} & \textbf{Bin width at 50 Hz}
    & \textbf{Effect} \\
    \midrule
    16 & 9  & 3.13 Hz & Few features, low resolution \\
    32 & 17 & 1.56 Hz & The setting of the kit lab \\
    64 & 33 & 0.78 Hz & Many features, high resolution \\
    \bottomrule
  \end{tabular}
  \vskip 0.4em
  \keypoint{The FFT length is a design decision: more bins separate
    frequencies better, and they make the model input larger.}
\end{frame}

\begin{frame}{Spectral power features}
  \begingroup\centering
  \begin{tikzpicture}[font=\scriptsize, x=0.078cm, y=0.8cm]
    % A window of 100 samples and four frames of 32 samples.
    \draw[coolgray, fill=boxgray] (0,1.6) rectangle (100,2.2);
    \node at (50,1.9) {window: 100 samples of one axis};
    \foreach \a/\b in {0/32, 32/64, 64/96} {
      \draw[darkcyan, fill=darkcyan!20] (\a,0.6) rectangle (\b,1.2);
    }
    \draw[darkcyan, fill=darkcyan!20] (96,0.6) rectangle (100,1.2);
    \draw[darkcyan, dashed] (100,0.6) rectangle (128,1.2);
    \node at (16,0.9) {frame 1};
    \node at (48,0.9) {frame 2};
    \node at (80,0.9) {frame 3};
    \node at (114,0.9) {zeros};
    \node[anchor=west] at (0,0.15) {FFT of each frame, power of each bin,
      then the maximum of each bin over the frames};
  \end{tikzpicture}\par\endgroup
  \vskip 0.2em
  \small
  \begin{itemize}
    \item The window has 100 samples. The FFT length is 32. So the block
          cuts the window into frames of 32 samples.
    \item The power of a bin is the square of its FFT value, divided by the
          FFT length.
    \item For each bin, the block keeps the largest power over the frames.
    \item The bin at 0 Hz is not a feature: the block removed the mean
          before. The result is 16 power values for each axis.
  \end{itemize}
  \source{Method of the spectral analysis block of Edge Impulse:
    \emph{Machine Learning Systems}, hardware kits, ``DSP Spectral
    Features''.}
\end{frame}

\begin{frame}{The feature block of the kit lab}
  \small
  \renewcommand{\arraystretch}{1.3}
  \begin{tabular}{@{}L{0.36\textwidth}L{0.38\textwidth}r@{}}
    \toprule
    \textbf{Group} & \textbf{Features for each axis} & \textbf{Number} \\
    \midrule
    Statistics in the time domain & RMS, skewness, kurtosis & 3 \\
    Statistics of the spectrum & Skewness and kurtosis of the power values
        & 2 \\
    Spectral power & One value for each bin from 1 to 16 & 16 \\
    \midrule
    Sum for one axis & & 21 \\
    Sum for three axes & & \textbf{63} \\
    \bottomrule
  \end{tabular}
  \vskip 0.7em
  \begin{itemize}
    \item The block first removes the mean of each axis in the window.
    \item The 63 features are the input of the classifier: 63 neurons in
          the input layer, then 20, 10, and 4.
    \item Edge Impulse Studio calculates the same features. You compare
          the two paths in Part D of the lab.
  \end{itemize}
  \source{\emph{Machine Learning Systems}, XIAOML Kit, ``Motion
    Classification and Anomaly Detection'', and ``DSP Spectral Features''.}
\end{frame}

\begin{frame}[fragile]{The features of one axis in Python}
\begin{lstlisting}[style=python]
import numpy as np

def spectral_power(x, nfft=32):
    """Largest power of each frequency bin over the frames of the window."""
    power = np.zeros(nfft // 2 + 1)
    for start in range(0, len(x), nfft):
        frame = x[start:start + nfft]              # the last frame is short
        spectrum = np.abs(np.fft.rfft(frame, nfft))  # rfft adds zeros
        power = np.maximum(power, spectrum ** 2 / nfft)
    return power[1:]                               # no bin at 0 Hz

x = window[:, 0]                # 100 samples of the x axis
x = x - x.mean()                # remove gravity and the offset
rms = np.sqrt(np.mean(x ** 2))
power = spectral_power(x)       # 16 values
\end{lstlisting}
  \small
  \begin{itemize}
    \item The training notebook calculates the features of each window with
          this code.
    \item The sketch on the board calculates the same features in C++. The
          two results must be equal for the same window.
  \end{itemize}
  \source{The function follows \texttt{welch\_max\_hold} of ``DSP Spectral
    Features'' in \emph{Machine Learning Systems}.}
\end{frame}

% ------------------------------------------------------------- smaller model
\begin{frame}{The model becomes smaller}
  \small
  \renewcommand{\arraystretch}{1.3}
  \begin{tabular}{@{}L{0.38\textwidth}rr@{}}
    \toprule
    & \textbf{Raw window} & \textbf{Features} \\
    \midrule
    Input values                        & 300   & 63 \\
    Parameters of the first layer (20 neurons) & 6020 & 1280 \\
    Parameters of the complete model    & 6274  & 1534 \\
    Weights in \code{float32}, in bytes & 25\,096 & 6136 \\
    Peak of the activations, in bytes   & 1280  & 332 \\
    \bottomrule
  \end{tabular}
  \vskip 0.7em
  \begin{itemize}
    \item The two models have the same hidden layers: 20 and 10 neurons,
          and 4 outputs.
    \item In a fully connected network, the first layer holds almost all
          parameters. Its size is the number of inputs times the number of
          neurons.
    \item The signal processing is not free. It needs code in the flash
          and time on the processor. You measure the two in the lab.
  \end{itemize}
\end{frame}

\begin{frame}{Why features need less model: invariance}
  \begingroup\centering
  \begin{tikzpicture}[font=\scriptsize, x=3.0cm, y=0.6cm]
    \foreach \o/\p/\t in {0/0/window 1, 1.25/110/window 2: 0.15 s later} {
      \draw[->, coolgray] (\o,0) -- (\o+1.08,0);
      \draw[darkcyan, line width=0.9pt]
        plot[domain=0:1, samples=100, smooth] (\x+\o,{sin(360*2*\x+\p)});
      \node[anchor=south] at (\o+0.5,1.2) {\t};
    }
    \node[align=center, anchor=west] at (2.55,0.3)
      {300 different\\numbers};
    \node[cardinalred, align=center, anchor=west] at (2.55,-1.3)
      {the same RMS,\\the same spectrum};
  \end{tikzpicture}\par\endgroup
  \vskip 0.4em
  \begin{itemize}
    \item The two windows hold the same motion. The second window starts
          later.
    \item As raw input, the two windows are two different vectors. The
          network must learn that they are the same class, for each
          possible shift.
    \item The RMS and the spectral power do not change with the shift. The
          features hold what the network must learn if not.
  \end{itemize}
  \vskip 0.2em
  \keypoint{A feature that removes a variation saves the parameters and
    the training data that the network needs to learn that variation.}
\end{frame}

\begin{frame}{The cost moves to the signal processing}
  \small
  \renewcommand{\arraystretch}{1.3}
  \begin{tabular}{@{}L{0.20\textwidth}L{0.34\textwidth}L{0.36\textwidth}@{}}
    \toprule
    & \textbf{Raw window} & \textbf{Features} \\
    \midrule
    Flash   & A larger model & A smaller model, and the code of the
                               features \\
    RAM     & A larger arena & A smaller arena, and buffers for the FFT \\
    Latency & The inference only & The features, then a shorter inference \\
    Design  & No knowledge of the signal is necessary
            & You must select the features \\
    Data    & More training data is necessary & Less training data is
                                                 necessary \\
    \bottomrule
  \end{tabular}
  \vskip 0.7em
  \begin{itemize}
    \item The pre-processing is a part of the latency. Measure the time of
          the features and the time of \code{Invoke()} separately.
    \item An FFT of 32 points is a small calculation. For a microphone at
          16 kHz, the signal processing can need more time than the model.
  \end{itemize}
\end{frame}

\begin{frame}{The same features in training and on the device}
  \begingroup\centering
  \begin{tikzpicture}[font=\footnotesize, >={Stealth[length=5pt]},
      b/.style={draw=coolgray, rounded corners=3pt, fill=boxgray,
                minimum width=2.4cm, minimum height=0.9cm, align=center}]
    \node[b] (w) at (0,-0.75) {One test\\window};
    \node[b] (p) at (3.6,0)    {Features in\\Python};
    \node[b] (c) at (3.6,-1.5) {Features in\\C++, on the board};
    \node[b, draw=cardinalred] (d) at (7.6,-0.75) {Compare the\\63 numbers};
    \draw[->] (w.east) -- (p.west);
    \draw[->] (w.east) -- (c.west);
    \draw[->] (p.east) -- (d.west);
    \draw[->] (c.east) -- (d.west);
  \end{tikzpicture}\par\endgroup
  \vskip 0.5em
  \begin{itemize}
    \item The network learned from the features of the Python code. On the
          board, C++ code calculates the features.
    \item A small difference changes the input of the model: a different
          FFT length, no mean removal, a different unit.
    \item This is the training-serving skew of Day 2. It gives no error
          message.
  \end{itemize}
  \vskip 0.2em
  \keypoint{Test the feature code of the board with a window that has a
    known result. Do this before you test the model.}
\end{frame}

\begin{frame}{Features by hand, or learned features}
  \small
  \renewcommand{\arraystretch}{1.3}
  \begin{tabular}{@{}L{0.17\textwidth}L{0.36\textwidth}L{0.37\textwidth}@{}}
    \toprule
    \textbf{Sensor} & \textbf{Usual input of the model} & \textbf{Reason} \\
    \midrule
    IMU        & Statistical and spectral features, with a small fully
                 connected network
               & The physics of the motion is known. The model is very
                 small. \\
    Microphone & A spectrogram or MFCC features, with a small convolutional
                 network (Day 5)
               & A raw second of audio has 16\,000 values. The spectrum
                 has much less. \\
    Camera     & The pixels, with a convolutional network (Day 5)
               & No simple formula describes an object. The network learns
                 its filters. \\
    \bottomrule
  \end{tabular}
  \vskip 0.7em
  \begin{itemize}
    \item A convolution layer on a raw time series is a set of filters that
          the training finds. It can replace the FFT, with more
          parameters and more data.
    \item When you know which property of the signal separates the classes,
          calculate that property directly.
  \end{itemize}
\end{frame}

% ------------------------------------------------------------------ lab brief
\begin{frame}{The lab: two paths to the same classifier}
  \begingroup\centering
  \begin{tikzpicture}[font=\footnotesize, >={Stealth[length=5pt]},
      b/.style={draw=coolgray, rounded corners=3pt, fill=boxgray,
                minimum width=1.95cm, minimum height=0.9cm, align=center}]
    \node[b] (d) at (0,-0.8) {Dataset of\\Day 2};
    \node[b] (a1) at (2.9,0)  {Notebook:\\features, Keras};
    \node[b] (a2) at (5.6,0)  {LiteRT file,\\C array};
    \node[b, draw=cardinalred] (a3) at (8.3,0) {Your sketch\\with TFLM};
    \node[b] (b1) at (2.9,-1.6) {Edge Impulse\\Studio};
    \node[b] (b2) at (5.6,-1.6) {Arduino\\library};
    \node[b, draw=cardinalred] (b3) at (8.3,-1.6) {Sketch of\\the kit lab};
    \draw[->] (d.east) -- (a1.west); \draw[->] (d.east) -- (b1.west);
    \draw[->] (a1) -- (a2); \draw[->] (a2) -- (a3);
    \draw[->] (b1) -- (b2); \draw[->] (b2) -- (b3);
  \end{tikzpicture}\par\endgroup
  \vskip 0.5em
  \small
  \begin{itemize}
    \item \textbf{Part A:} calculate the features of your windows and train
          a small classifier.
    \item \textbf{Part B:} convert the model, write the inference code, and
          show the class on the display.
    \item \textbf{Part C:} measure the arena, the flash, and the latency,
          with PSRAM and with no PSRAM.
    \item \textbf{Part D:} train the same task in Edge Impulse Studio, and
          compare the two results in one table.
  \end{itemize}
  \keypoint{The two paths use the same data and the same type of features.
    The difference is who writes the code, and which kernels run.}
\end{frame}

% -------------------------------------------------------------------- exercise
\exerciseframe{raw windows and spectral features}{%
  A motion classifier uses a window of 2 s with 3 axes. The first layer of
  the network has 20 neurons. The feature block gives 21 features for each
  axis.
  \begin{enumerate}
    \item The sampling rate is 50 Hz. How many values has the raw window?
          How many features has the feature block?
    \item Calculate the parameters of the first layer for the two inputs.
    \item The team changes the sampling rate to 100 Hz. The window stays at
          2 s, and the FFT length stays at 32. Calculate the two input
          sizes and the parameters of the first layer again.
    \item At 100 Hz with an FFT length of 32, what is the width of one
          frequency bin? What is the highest frequency of the spectrum?
    \item Which input do you select for a device with 2 KB of free RAM?
          Give the reason.
  \end{enumerate}
  Work with your neighbour for 8 minutes.
}

\answerframe{raw windows and spectral features}{%
  \small
  \begin{enumerate}
    \item Raw window: $2 \times 50 \times 3 = 300$ values. Features:
          $21 \times 3 = 63$.
    \item Raw window: $300 \times 20 + 20 = 6020$. Features:
          $63 \times 20 + 20 = 1280$.
    \item Raw window: $2 \times 100 \times 3 = 600$ values and
          $600 \times 20 + 20 = 12\,020$ parameters. Features: 63 values
          and 1280 parameters, with no change. The number of features
          depends on the FFT length, not on the sampling rate.
    \item Bin width: $100 / 32 = 3.13$ Hz. Highest frequency: 50 Hz. The
          same number of bins now covers two times the frequency range, so
          each bin is two times as wide.
    \item The features. With \code{float32}, the raw window at 50 Hz needs
          1200 bytes for the input and 80 bytes for the first output: 1280
          bytes of the 2 KB, with no space for the data of the
          interpreter. The features need $63 \times 4 + 80 = 332$ bytes.
  \end{enumerate}
}
